\documentclass[10pt]{ctexart}
\usepackage[a4paper,margin=0.9in]{geometry}
\usepackage{amsmath,amssymb,booktabs,array}
\usepackage{enumitem}
\usepackage{listings}
\usepackage{xcolor}
\usepackage{microtype}
\usepackage{hyperref}
\setlist[itemize]{leftmargin=1.5em}
\setlist[enumerate]{leftmargin=1.8em}
\lstset{basicstyle=\ttfamily\small,breaklines=true,columns=fullflexible,keepspaces=true,frame=single,framerule=0.2pt,xleftmargin=0.5em,xrightmargin=0.5em}

\title{019：Soft-GiGPO：基于模型既视感的 Step 软匹配设想}
\author{}
\date{2026-06-10}

\begin{document}
\maketitle
\vspace{-1.2em}

\section*{摘要}
本文草拟一个新的 GiGPO 扩展方向：把当前基于 anchor observation 的硬文本匹配，替换或增强为基于模型“既视感”的 soft semantic matching。核心思想是：两个 step 不必完全文本相同才可以互相作为 step-level baseline；只要它们在模型表示空间、文本 embedding 空间或 frozen actor hidden space 中足够相似，就可以按连续权重参与对方的 relative advantage 估计。

这个方向暂命名为 \textbf{Soft-GiGPO / Semantic Step Matching}。它不改变当前 StepPPO-v3 的 value-head-only 定义，也不默认引入 action-level PPO ratio 或新的 critic。第一版建议只替换 GiGPO 的 step relative advantage estimator，保持 rollout、policy loss 和训练 recipe 不变。

\section{背景：当前 GiGPO 的硬匹配}
当前 GiGPO 的关键设计是：在同一个任务 \path{uid} 下，把处在相同 anchor state 的 rollout step 放入同一个 step group，然后在 group 内计算相对 step advantage。这个设计比只看 episode-level reward 更细，因为它试图回答：在“相同状态”下，不同 response/action 的未来 return 谁更好。

代码路径大致是：
\begin{lstlisting}
gigpo/core_gigpo.py
  compute_gigpo_outcome_advantage(...)
    step_group_uids = build_step_group(anchor_obs, index, enable_similarity, similarity_thresh)
    step_advantages = step_norm_reward(step_rewards, response_mask, step_group_uids, ...)
\end{lstlisting}

当前实现仍然主要是硬分组：默认 exact match；可选 \path{enable_similarity=True} 时使用 \path{SequenceMatcher} 文本相似度和阈值做 hard clustering。一旦两个 step 没落入同一 cluster，它们对彼此的 baseline/advantage 就完全没有贡献。

这种硬匹配在 WebShop 这样的文本环境中有明显限制：页面 observation 可能只因为商品排序、空格、价格显示、历史上下文或轻微 wording 改变而无法 exact match；反过来，两个文本相似的页面也可能在任务语义上不同。于是 step-level signal 可能既稀疏又脆弱。

\section{直觉：模型的“既视感”}
用户提出的“模型既视感”可以理解为：当模型看到当前 step context 时，它内部可能已经知道这个状态和 batch 内哪些历史或并行 rollout state 很像。这个像不像不一定等价于文本完全相同，而可能包含：
\begin{itemize}
  \item 当前商品搜索页面是否语义相近；
  \item 当前 observation 是否处在相似的搜索、筛选或详情页阶段；
  \item 当前 action space 是否类似；
  \item 当前任务进展是否类似；
  \item 当前 hidden representation 是否显示出相似的决策处境。
\end{itemize}

如果这种相似度可用，就可以把 GiGPO 的“同 anchor state 内相对比较”推广成“相似 state 之间的软相对比较”。一个 step 不再只和 exact-matched peers 比较，而是和一组相似 step 按权重比较。

\section{方法定义}
设一个 batch 中有 step 样本 $i=1,\dots,N$。每个样本属于任务组 $u_i$，有 anchor/context 表示 $s_i$，step-level return 或 discounted step reward 为 $R_i$。原始 GiGPO 会构造 hard group $G_i$，然后计算：
\begin{equation}
  A_i^{\mathrm{hard}} = R_i - \frac{1}{|G_i|}\sum_{j\in G_i} R_j,
\end{equation}
或在 \path{mean_std_norm} 下再除以 group std。

Soft-GiGPO 改为先定义相似表示和相似度：
\begin{equation}
  z_i=f_{\mathrm{match}}(s_i),\qquad
  \mathrm{sim}_{ij}=\cos(z_i,z_j).
\end{equation}
其中 $f_{\mathrm{match}}$ 可以是 frozen embedding model、actor selected hidden state、value head 的中间表示，或一个单独训练的轻量 projection。

在同一 task uid 内计算 kernel 权重：
\begin{equation}
  w_{ij}=\frac{\mathbf{1}[u_i=u_j]\mathbf{1}[i\ne j]\exp(\mathrm{sim}_{ij}/\tau)}
  {\sum_{k:u_k=u_i,k\ne i}\exp(\mathrm{sim}_{ik}/\tau)+\epsilon}.
\end{equation}
这里 $\tau$ 是 temperature，控制 soft matching 的尖锐程度。若 $\tau$ 很小，方法接近 nearest-neighbor；若 $\tau$ 较大，则接近 uid 内平均 baseline。

soft baseline 为：
\begin{equation}
  b_i^{\mathrm{soft}}=\sum_{j:u_j=u_i,j\ne i} w_{ij}R_j.
\end{equation}
soft step advantage 为：
\begin{equation}
  A_i^{\mathrm{soft}}=R_i-b_i^{\mathrm{soft}}.
\end{equation}

如果需要标准化，可以用加权方差：
\begin{equation}
  \sigma_i^2=\sum_j w_{ij}(R_j-b_i^{\mathrm{soft}})^2,
  \qquad
  \hat A_i^{\mathrm{soft}}=\frac{R_i-b_i^{\mathrm{soft}}}{\sqrt{\sigma_i^2+\epsilon}}.
\end{equation}
为了避免自身 reward 泄漏，默认应使用 leave-one-out，即 $w_{ii}=0$。

\section{和 GiGPO 的关系}
Soft-GiGPO 不是推翻 GiGPO，而是把 hard anchor group 泛化成 soft anchor neighborhood。原始 GiGPO 可以看成一种特殊 kernel：
\begin{equation}
  w_{ij}^{\mathrm{hard}}\propto \mathbf{1}[g_i=g_j]\mathbf{1}[i\ne j].
\end{equation}
Soft-GiGPO 则把 indicator 变成连续 similarity kernel：
\begin{equation}
  w_{ij}^{\mathrm{soft}}\propto K(s_i,s_j).
\end{equation}
因此它保留 GiGPO 的核心精神：仍然是 group-relative / state-relative advantage；区别只是“状态相同”的定义从离散硬匹配变成连续软匹配。

\section{“模型既视感”的实现选择}
\subsection{文本 embedding 版}
对 anchor observation 或 pre-action context 做 embedding：
\begin{equation}
  z_i=\mathrm{Embed}(\mathrm{text}(s_i)).
\end{equation}
优点是实现快，不依赖 actor forward 额外改动；缺点是外部 embedding 不一定理解 WebShop action semantics，也可能引入额外模型和 runtime。

\subsection{Actor hidden state 版}
使用 actor 在准备生成 action 前的 selected hidden state：
\begin{equation}
  z_i=h_{\mathrm{pre\text{-}action\text{-}last\text{-}context\text{-}token},i}.
\end{equation}
这和当前 v3 value head 使用的 state representation 接近。优点是“既视感”来自同一个 policy model，可能更贴近策略决策；缺点是 representation 非平稳，训练过程中会漂移。

\subsection{Frozen actor snapshot 版}
为避免 representation non-stationarity，可以用 frozen checkpoint 或 reference actor 计算 $z_i$，例如初始 SFT actor、GiGPO baseline checkpoint，或周期性冻结的 target encoder。优点是相似度稳定；缺点是可能逐渐落后于当前 policy 的状态分布。

\subsection{Value-aware 版}
如果 v3 value head 已证明 selected hidden state 有 step return signal，可以用 value prediction 辅助相似度：
\begin{equation}
  \mathrm{sim}_{ij}=\alpha\cos(z_i,z_j)-\beta |V(s_i)-V(s_j)|.
\end{equation}
直觉是：两个状态不仅文本或 representation 相似，而且预测 return 也相近，才更适合作为 baseline peer。这个版本风险较高，因为它会把 learned value 的 bias 引入 advantage estimator，应作为后续 ablation。

\section{推荐第一版：Soft-GiGPO-v0}
第一版建议只替换 GiGPO 的 step group / step advantage 计算，不改 policy loss，不接入 StepPPO-v3 value head，不改变 rollout recipe。

\begin{lstlisting}
algorithm:
  adv_estimator: gigpo
  gigpo:
    step_advantage_w: 1.0
    mode: mean_norm
    soft_matching:
      enable: true
      encoder: text_embedding        # text_embedding / actor_hidden / frozen_actor
      scope: same_uid                # only match within the same task uid
      kernel: cosine_softmax
      temperature: 0.1
      top_k: 8
      min_neighbors: 2
      leave_one_out: true
      fallback: hard_or_uid_mean
      detach_embeddings: true
\end{lstlisting}

计算流程：
\begin{enumerate}
  \item 对每个 step 构造匹配表示 $z_i$。
  \item 在同一 \path{uid} 内计算 pairwise similarity。
  \item 取 top-k neighbors，做 softmax 权重。
  \item 用 weighted leave-one-out return baseline 得到 $A_i^{\mathrm{soft}}$。
  \item 把 scalar step advantage broadcast 到 response tokens，保持原 GiGPO actor update 不变。
\end{enumerate}

\section{关键设计选择}
\paragraph{匹配范围限制在同一任务 uid 内。}不同 WebShop task 的目标商品不同，即使页面文本相似，reward 语义也可能完全不同。第一版应只允许同一 \path{uid} 内的 rollout 相互匹配。

\paragraph{使用 top-k 而不是全连接 dense kernel。}全连接 softmax 可能把大量弱相关样本纳入 baseline，导致 advantage 过度平滑。top-k 能保留“相似 neighborhood”的直觉，并降低计算开销。

\paragraph{默认 leave-one-out。}如果把自己放进 baseline，尤其当 top-k 很小或 similarity 很尖锐时，会把 $R_i$ 泄漏进 $b_i$，使 advantage 被人为压小。因此默认 $w_{ii}=0$。

\paragraph{fallback 很重要。}如果某个 uid 内有效 neighbor 太少，或相似度全都很低，应该 fallback 到原始 hard GiGPO group、uid 内 mean baseline，或直接把 step advantage 置零，只保留 episode advantage。第一版建议 \path{fallback=hard_or_uid_mean}。

\section{诊断指标}
Soft matching 是否有用，不能只看 validation score。必须先确认 matching 质量和 advantage 行为：
\begin{lstlisting}
gigpo_soft/neighbor_count_mean
gigpo_soft/effective_neighbors_mean
gigpo_soft/top1_similarity_mean
gigpo_soft/topk_similarity_mean
gigpo_soft/weight_entropy_mean
gigpo_soft/fallback_ratio
gigpo_soft/soft_baseline_mean
gigpo_soft/soft_adv_mean
gigpo_soft/soft_adv_std
gigpo_soft/hard_soft_adv_corr
gigpo_soft/group_size_equiv_mean
\end{lstlisting}

其中 effective neighbors 可以定义为：
\begin{equation}
  N_{\mathrm{eff},i}=\frac{1}{\sum_j w_{ij}^2}.
\end{equation}
它能判断 soft kernel 是退化成单邻居 hard match，还是退化成过宽平均。

还应做 offline diagnostic：hard group size vs soft effective neighbor size；hard advantage 与 soft advantage 的相关性；top-k neighbors 的 observation 文本人工抽样；soft similarity 与 reward gap 的关系；soft matching 是否降低 singleton 比例；soft advantage 是否减少极端 outlier。

\section{可能收益与主要风险}
\subsection{可能收益}
\begin{enumerate}
  \item 缓解 exact-match 稀疏性：相似但不完全相同的 states 可以互相提供 baseline。
  \item 更细粒度 credit assignment：不必等到完全同一 anchor observation 才比较 response/action。
  \item 提高 sample efficiency：同一 batch 内更多 step 能参与 step-relative normalization。
  \item 兼容当前 GiGPO pipeline：第一版只替换 step advantage，不改 actor loss、ratio、rollout。
  \item 为 value head 提供诊断桥梁：如果 actor hidden state 的 soft matching 有用，说明 representation 中已有可用 state similarity signal。
\end{enumerate}

\subsection{主要风险}
\begin{enumerate}
  \item 错误匹配引入 bias：语义相似不等于决策等价，错误 neighbors 会污染 baseline。
  \item 过度平滑 advantage：soft baseline 过宽会把有用差异抹掉。
  \item representation 漂移：若用当前 actor hidden state，similarity 空间会随训练变化。
  \item reward leakage：不做 leave-one-out 会低估 advantage。
  \item 额外计算开销：pairwise similarity 需要 uid 内 top-k 和缓存。
  \item 与 v3 value head 混淆：当前 StepPPO-v3 主线仍应保持 value-head-only；soft matching 应作为新的 GiGPO ablation 或未来版本。
\end{enumerate}

\section{实验计划}
\subsection{Offline matching study}
先不训练，只用已有 GiGPO rollout/log dump 或新增少量 rollout dump：
\begin{enumerate}
  \item 对每个 batch 构造 hard groups 和 soft neighborhoods。
  \item 统计 singleton hard groups 被 soft matching 扩展后的 effective neighbor 数。
  \item 人工检查 top-k neighbors 是否语义合理。
  \item 比较 hard advantage 与 soft advantage 的均值、方差、极端值和相关性。
  \item 检查 soft similarity 是否和 reward gap / future return gap 有负相关。
\end{enumerate}
如果 soft neighbors 看起来不合理，不应直接进入 full training。

\subsection{2GPU smoke}
目标不是追分，而是看 estimator 是否稳定：valid action ratio 不快速崩；response clip ratio 不快速升高；soft fallback ratio 不长期接近 1；effective neighbors 不塌缩到 1，也不过宽接近 uid group size；policy loss、KL、grad norm 不异常。

\subsection{4GPU seed 2026 full run}
只和已有 GiGPO seed 2026 对齐：同 batch size、rollout group size、test freq、GPU 数。比较 validation task score / success、valid action ratio、response length / clip ratio、runtime、hard-soft advantage diagnostics。

\subsection{Ablation 顺序}
\begin{center}
\begin{tabular}{lll}
\toprule
实验 & 改动 & 目的 \\
\midrule
GiGPO hard & 当前 baseline & 对照 \\
Soft-GiGPO text-emb top-k & 外部/frozen text embedding & 验证软匹配基本收益 \\
Soft-GiGPO actor-hidden frozen & frozen actor representation & 验证模型既视感是否优于文本 embedding \\
Soft-GiGPO actor-hidden online & 当前 actor hidden，detach & 验证 online representation 是否可用 \\
Soft-GiGPO hard+soft mix & $A=A_{hard}+\lambda A_{soft}$ & 降低纯 soft 风险 \\
\bottomrule
\end{tabular}
\end{center}

第一轮不建议同时加入 value head、action-level PPO ratio 或 StepPPO-v3 其它改动。

\section{代码实现草图}
可以新增：
\begin{lstlisting}
def compute_soft_step_advantage(
    step_rewards: torch.Tensor,
    response_mask: torch.Tensor,
    uid: np.ndarray,
    match_embeddings: torch.Tensor,
    temperature: float = 0.1,
    top_k: int = 8,
    remove_std: bool = True,
    leave_one_out: bool = True,
    epsilon: float = 1e-6,
):
    # 1. split rows by uid
    # 2. cosine similarity within each uid
    # 3. mask self, take top-k
    # 4. softmax(sim / temperature)
    # 5. weighted baseline and optional weighted std
    # 6. broadcast scalar advantage to response tokens
    return step_advantages, diagnostics
\end{lstlisting}

并在 \path{compute_gigpo_outcome_advantage} 中加开关：
\begin{lstlisting}
if soft_matching.enable:
    step_advantages, soft_metrics = compute_soft_step_advantage(...)
else:
    step_group_uids = build_step_group(...)
    step_advantages = step_norm_reward(...)
\end{lstlisting}

如果第一版使用 text embedding，embedding 可以先在 trainer driver 侧生成并放入 \path{DataProto.batch["step_match_embeddings"]}。如果使用 actor hidden state，则可以复用 v3 value head 的 selected-state extraction 思路，但第一版应只 dump/compute embedding，不训练 value head。

\section{与 StepPPO-v3 的边界}
这份设想不改变当前 StepPPO-v3 定义。当前 v3 仍然是：
\begin{lstlisting}
GiGPO policy training + actor-side shared value head
\end{lstlisting}

Soft-GiGPO 是另一个方向：改 GiGPO 的 step relative advantage estimator，把 hard anchor matching 换成 soft semantic matching。它可以作为 GiGPO ablation、后续 GiGPO-v2 / Soft-GiGPO，或将来 StepPPO-v4 的一个可能组件。但在当前阶段，不应把 soft matching、value head、action-level PPO ratio 一次性合并，否则很难定位收益或退化来源。

\section{初步结论}
这个想法值得做，而且比直接把 value head 接入 policy advantage 更稳健：它仍然保留 GiGPO 的 relative baseline 思路，只是把“相同 state”的硬文本定义放宽成“相似 state”的软语义定义。

建议下一步先做 offline matching study：用已有 GiGPO 或 v3 rollout batch，比较 hard group 与 soft neighborhood 的质量。如果模型既视感能稳定找到语义合理、return 分布相近的 peers，再进入 2GPU smoke 和 4GPU full run。

\end{document}
